% LuaLaTeX用の設定
\documentclass[a4paper,11pt]{ltjsarticle}

% --- パッケージ ---
\usepackage{graphicx}
\usepackage{amsmath,amssymb}
\usepackage{bm}
\usepackage{url}

% --- タイトル情報 ---
\title{第1回 レポート課題}
\author{氏名：安藤 太郎 \\ 学籍番号：XXXXXXXX}
\date{\today}

\begin{document}

\maketitle


\section{はじめに}
このレポートでは、参考文献を外部ファイルを使わずに直接本文に記述しています。
数個の引用であれば、この方式が最もシンプルでトラブルが少ないです。

\section{引用の方法}
本文中で引用するには、このように書きます\cite{book1}。
複数の文献がある場合も同様です\cite{paper1}。

\section{おわりに}
レポートの最後に `thebibliography` 環境を作ることで、自動的に「参考文献」という見出しとともにリストが表示されます。

% --- 参考文献（直接入力方式） ---
% {99} は「最大99個まで書く」というおまじないです。
\begin{thebibliography}{99}

    \bibitem{book1} 
    安藤 崇央, 『ほげほげ学入門』, ほげほげ評論社, 2026.

    \bibitem{paper1} 
    長島 史苑, 「ふがふがの研究について」, ふがふが学会論文誌, Vol.10, pp.1-5, 2024.

\end{thebibliography}

\end{document}